Dans un lycée, 120 élèves sont sportifs de haut niveau et sont séparés en deux sections : la section aquatique et la section judo.

	Pour chaque section, les élèves peuvent être en seconde, en première ou en terminale.

	La répartition de ces 120 élèves est donnée dans le tableau ci-dessous.

	\begin{center}
		\begin{tblr}{colspec={l *{3}{X[c,2cm]}X[c,1cm]},
				hline{1}={2-Z}{solid},hline{2-Z}={solid},
				vline{1}={2-Z}{solid},vline{2-Z}={solid}}
			                  & Seconde & Première & Terminale & Total \\
			Section judo      &         &          & 8         & 24    \\
			Section aquatique & 40      &          &           &       \\
			Total             & 50      &          & 14        & 120
		\end{tblr}
	\end{center}

On choisit au hasard un élève parmi les sportifs de haut niveau.

On considère les évènements suivants :

\begin{multicols}{2}
	\begin{itemize}[label={\textbullet~}]
	\item $S$ : «\,l'élève est en seconde\,»
	\item $P$ : «\,l'élève est en première\,»
	\item $T$ : «\,l'élève est en terminale\,»
	\item $A$ : «\,l'élève est en section aquatique\,»
	\item $J$ : «\,l'élève est en section judo\,»
	\end{itemize}

	~~
	\end{multicols}

	Pour un évènement quelconque $E$ : on note $P(E)$ la probabilité de $E$ et $P_{F}(E)$ la probabilité de $E$ sachant l'évènement $F$.

	\begin{enumerate}
		\item Reproduire le tableau ci-dessus en complétant les données manquantes.
	\end{enumerate}

	\emph{Dans les questions qui suivent, les résultats seront donnés sous forme de fraction.}

	\begin{enumerate}[resume]
		\item Décrire l'évènement $A \cap S$ à l'aide d'une phrase puis calculer sa probabilité.

		\item On choisit au hasard un élève sportif en seconde, calculer la probabilité qu'il soit en section aquatique.

		\item \begin{enumerate}
			\item Calculer la probabilité de l'évènement $J$.
			\item Calculer $\mathrm{P}_{T}(J)$.
			\item Les évènements\, $J$ et $T$ sont-ils indépendants ?

			Justifier votre réponse.
		\end{enumerate}
	\end{enumerate}
